% Personal statement for the Kaiser Permanente Division of Research Delivery Science Fellowship
% (July 2027 cohort, ORCHESTRATE-HF). The form asks for at most 2-3 pages covering prior record,
% reasons for interest in the fellowship and plans for the future.
% Shares cv_preamble.tex and cv_header.tex with the CVs so the documents match.
\documentclass[11pt, letterpaper]{article}
% Shared preamble for main.tex (general CV) and cv_health_research.tex
% Packages:
\usepackage[
ignoreheadfoot, % set margins without considering header and footer
top=2 cm, % seperation between body and page edge from the top
bottom=2 cm, % seperation between body and page edge from the bottom
left=2 cm, % seperation between body and page edge from the left
right=2 cm, % seperation between body and page edge from the right
footskip=1.0 cm, % seperation between body and footer
% showframe % for debugging
]{geometry} % for adjusting page geometry
\usepackage{titlesec} % for customizing section titles
\usepackage{tabularx} % for making tables with fixed width columns
\usepackage{array} % tabularx requires this
\usepackage[dvipsnames]{xcolor} % for coloring text
\definecolor{primaryColor}{RGB}{0, 79, 144} % define primary color
\usepackage{enumitem} % for customizing lists
\usepackage{fontawesome5} % for using icons
\usepackage{amsmath} % for math
\usepackage[
pdftitle={Bipin Lekhak's CV},
pdfauthor={Bipin Lekhak},
pdfcreator={LaTeX with RenderCV},
colorlinks=true,
urlcolor=primaryColor
]{hyperref} % for links, metadata and bookmarks
\usepackage[pscoord]{eso-pic} % for floating text on the page
\usepackage{calc} % for calculating lengths
\usepackage{bookmark} % for bookmarks
\usepackage{lastpage} % for getting the total number of pages
\usepackage{changepage} % for one column entries (adjustwidth environment)
\usepackage{paracol} % for two and three column entries
\usepackage{ifthen} % for conditional statements
\usepackage{needspace} % for avoiding page brake right after the section title
\usepackage{iftex} % check if engine is pdflatex, xetex or luatex
\usepackage{etoolbox} % for \patchcmd and \AtBeginEnvironment

% Ensure that generate pdf is machine readable/ATS parsable:
\ifPDFTeX
    \input{glyphtounicode}
    \pdfgentounicode=1
    % \usepackage[T1]{fontenc} % this breaks sb2nov
    \usepackage[utf8]{inputenc}
    \usepackage{lmodern}
\fi



% Some settings:
\AtBeginEnvironment{adjustwidth}{\partopsep0pt} % remove space before adjustwidth environment
\pagestyle{empty} % no header or footer
\setcounter{secnumdepth}{0} % no section numbering
\setlength{\parindent}{0pt} % no indentation
\setlength{\topskip}{0pt} % no top skip
\setlength{\columnsep}{0cm} % set column seperation
\makeatletter
\let\ps@customFooterStyle\ps@plain % Copy the plain style to customFooterStyle
\patchcmd{\ps@customFooterStyle}{\thepage}{
\color{gray}\textit{\small Bipin Lekhak - Page \thepage{} of \pageref*{LastPage}}
}{}{} % replace number by desired string
\makeatother
\pagestyle{customFooterStyle}

\titleformat{\section}{\needspace{4\baselineskip}\bfseries\large\color{primaryColor}}{}{0pt}{}[\vspace{1pt}\titlerule]

\titlespacing{\section}{
% left space:
-1pt
}{
% top space:
0.3 cm
}{
% bottom space:
0.2 cm
} % section title spacing

\renewcommand\labelitemi{$\circ$} % custom bullet points
\newenvironment{highlights}{
\begin{itemize}[
    topsep=0.10 cm,
    parsep=0.10 cm,
    partopsep=0pt,
    itemsep=0pt,
    leftmargin=0.4 cm + 10pt
    ]
    }{
\end{itemize}
} % new environment for highlights

\newenvironment{highlightsforbulletentries}{
\begin{itemize}[
    topsep=0.10 cm,
    parsep=0.10 cm,
    partopsep=0pt,
    itemsep=0pt,
    leftmargin=10pt
    ]
    }{
\end{itemize}
} % new environment for highlights for bullet entries


\newenvironment{onecolentry}{
\begin{adjustwidth}{
    0.2 cm + 0.00001 cm
    }{
    0.2 cm + 0.00001 cm
    }
    }{
\end{adjustwidth}
} % new environment for one column entries

\newenvironment{twocolentry}[2][]{
\onecolentry
\def\secondColumn{#2}
\setcolumnwidth{\fill, 4.5 cm}
\begin{paracol}{2}
    }{
    \switchcolumn \raggedleft \secondColumn
\end{paracol}
\endonecolentry
} % new environment for two column entries

\newenvironment{header}{
\setlength{\topsep}{0pt}\par\kern\topsep\centering\linespread{1.5}
}{
\par\kern\topsep
} % new environment for the header

\newcommand{\placelastupdatedtext}{ % \placetextbox{<horizontal pos>}{<vertical pos>}{<stuff>}
\AddToShipoutPictureFG*{ % Add <stuff> to current page foreground
\put(
\LenToUnit{\paperwidth-2 cm-0.2 cm+0.05cm},
\LenToUnit{\paperheight-1.0 cm}
){\vtop{{\null}\makebox[0pt][c]{\small\color{gray}\textit{Last updated in October 2026}\hspace{\widthof{Last updated in October 2026}}
}}} %
} %
} %

% save the original href command in a new command:
\let\hrefWithoutArrow\href

% new command for external links:
\renewcommand{\href}[2]{\hrefWithoutArrow{#1}{\ifthenelse{\equal{#2}{}}{ }{#2 }\raisebox{.15ex}{\footnotesize \faExternalLink*}}}

% one line in the skills section: bold label, then the list
\newcommand{\skillline}[2]{\textbf{#1:} #2}



\hypersetup{pdftitle={Bipin Lekhak - Personal Statement}}
\setlength{\parskip}{0.5em}

\begin{document}
    \renewcommand{\placelastupdatedtext}{} % no "last updated" stamp on the statement
    % Shared name and contact header, \input right after \begin{document}
    \newcommand{\AND}{\unskip
    \cleaders\copy\ANDbox\hskip\wd\ANDbox
    \ignorespaces
    }
    \newsavebox\ANDbox
    \sbox\ANDbox{}

    \placelastupdatedtext
    \begin{header}
        \textbf{\fontsize{24 pt}{24 pt}\selectfont Bipin Lekhak}

        \vspace{0.3 cm}

        \normalsize
        \mbox{{\color{black}\footnotesize\faMapMarker*}\hspace*{0.13cm}Rochester, NY} %
        \kern 0.25 cm %
        \AND %
        \kern 0.25 cm %
        \mbox{\hrefWithoutArrow{mailto:bl1282@rit.edu}{\color{black}{\footnotesize\faEnvelope[regular]}\hspace*{0.13cm}bl1282@rit.edu}} %
        \kern 0.25 cm %
        \AND %
        \kern 0.25 cm %
        \mbox{\hrefWithoutArrow{https://bipinlekhak.com.np/}{\color{black}{\footnotesize\faLink}\hspace*{0.13cm}bipinlekhak.com.np}}%

        \mbox{\hrefWithoutArrow{https://scholar.google.com/citations?user=z8T6GewAAAAJ}{\color{black}{\footnotesize\faGraduationCap}\hspace*{0.13cm}Google Scholar}}%
        \kern 0.25 cm%
        \AND%
        \kern 0.25 cm%
        \mbox{\hrefWithoutArrow{https://openreview.net/profile?id=~Bipin_Lekhak2}{\color{black}{\footnotesize\faFile*}\hspace*{0.13cm}OpenReview}}%
        \kern 0.25 cm%
        \AND%
        \kern 0.25 cm%
        \mbox{\hrefWithoutArrow{https://www.linkedin.com/in/bipin-lekhak/}{\color{black}{\footnotesize\faLinkedinIn}\hspace*{0.13cm}bipin-lekhak}}%
        \kern 0.25 cm%
        \AND%
        \kern 0.25 cm%
        \mbox{\hrefWithoutArrow{https://github.com/bipin-lekhak}{\color{black}{\footnotesize\faGithub}\hspace*{0.13cm}bipin-lekhak}} %
    \end{header}


    \vspace{0.2 cm}
    \begin{center}
        {\Large\bfseries\color{primaryColor} Personal Statement}

        \vspace{0.1 cm}
        Delivery Science Fellowship, Kaiser Permanente Division of Research
    \end{center}


    \section{From home monitoring to evaluating AI in care}

    In an NIH-funded study with the University of Rochester Medical Center (R01NR018301), heart-failure patients
    used a sensor-equipped toilet seat that recorded cardiac measurements at home for 90 days after discharge. I
    was the lead student for RIT's part of the study. I worked in the field on data collection and participant
    surveys, led the cleaning, validation and processing of the data, and contributed to the predictive models
    (logistic regression, decision trees and random forests). I was first author of the resulting paper
    (Computing in Cardiology 2023), and a fuller predictive analysis is being completed.

    Working in the field showed me where home monitoring breaks down: many patients who consented were never
    enrolled, often because the seat did not fit their toilet, and blood pressure was recorded far less
    consistently than heart rate. Which patients and measurements go missing decides what a model built on such
    data can learn. Ours could still flag patients at higher risk of hospitalisation from a week of seat data.
    ORCHESTRATE-HF moves from predicting risk to suggesting treatment, and evaluating those suggestions means
    estimating outcomes under treatment sequences patients did not receive. That is the subject of my
    dissertation.


    \section{Methods for treatments that change over time}

    My dissertation, \textit{Learning Dynamics for Forecasting Future Trajectories under External Interventions}
    (proposal defended July 2026), predicts a patient's trajectory under a treatment plan they have not yet
    received. Its main result, iASIDE (first author, under review), learns a patient's untreated disease course
    and their response to treatment as separate components, adapts both to each patient from their own history,
    and states when these dynamics are identifiable from observational data. On MIMIC-III and on a tumour-growth
    simulator with known counterfactual outcomes, I benchmarked it against marginal structural models with
    inverse probability weighting, g-computation, and machine-learning methods for time-varying treatments; its
    advantage was largest at long horizons and in heterogeneous patient populations. With Suraj Poudel, I am now
    adapting X- and R-learners to estimate individual treatment effects from longitudinal data.


    \section{Proposed work on ORCHESTRATE-HF}

    The agents will help clinicians use the several medications recommended for heart failure. Each medication
    decision depends on blood pressure, kidney function and potassium, which earlier decisions changed. That feedback is time-varying confounding, and it shapes each phase of the project.

    \begin{enumerate}[leftmargin=0.6 cm, topsep=0pt, itemsep=0.3em, parsep=0pt]
        \item \textbf{Evaluating agents on past KPNC records.} Scoring an agent's recommendations requires the
        outcomes patients would have had if clinicians had followed them. I would estimate these in target-trial
        emulations of medication strategies, with both g-computation and counterfactual trajectory models, and
        treat disagreement between the two as a warning before an agent reaches patients.
        \item \textbf{The trial's statistical analysis plan.} Randomisation gives the average effect of
        AI-assisted care. I would work on the secondary analyses: per-protocol effects that account for whether
        clinicians followed each suggestion, and heterogeneous effects, including by race and ethnicity, to test
        whether the benefit is shared equitably.
        \item \textbf{Monitoring agents after deployment.} Once agents change treatment, later records reflect
        their actions, and a risk model retrained on those records can learn that high-risk patients are low risk
        because they were treated. Separating the untreated course from the treatment response, the core of
        iASIDE, addresses this directly.
    \end{enumerate}


    \section{What I bring and what I need}

    I bring longitudinal causal methods, experience collecting and validating data from patients at home,
    reproducible software for large-scale experiments (a benchmark of more than ten methods run on HPC
    clusters), and four years as a machine-learning engineer before the PhD. I have
    not designed a randomised trial, written a grant or published in clinical journals; working with Drs.\ Ankeet
    Bhatt, Alan Go and Vincent Liu on ORCHESTRATE-HF would teach me all three. My goal is an independent research
    career in a learning health system, evaluating the AI tools it deploys. I expect to complete my PhD in June
    2027 and would join the July 2027 cohort.

\end{document}
